\documentclass{article}
\usepackage[UTF8]{ctex}
\usepackage{amsmath,amssymb,amsfonts}
\usepackage{enumitem}
\usepackage{geometry}
\geometry{a4paper, margin=1in}
% === 错题本排版优化（不参与题目标识） ===
\setlength{\parindent}{0pt}
\setlength{\parskip}{0.35em plus 0.10em minus 0.08em}
\setlength{\emergencystretch}{3em}
\setlist{itemsep=0.25em, topsep=0.40em, parsep=0pt}
\allowdisplaybreaks[2]
% === 错题本排版优化结束 ===
\begin{document}
\section*{西北大学2026年数学分析真题}
\section*{178．西北大学2026年数学分析试卷（回忆版）}

\subsection{题目1}
一．（15 分）解答如下问题：

\begin{itemize}
  \item[（1）] 用极限的定义证明：当 \(\lim _{n \rightarrow \infty} a_{n}=0\) 时，有 \(\lim _{n \rightarrow \infty} \frac{a_{1}+a_{2}+\cdots+a_{n}}{n}=0\) ．
  \item[（2）] 若 \(\lim _{n \rightarrow \infty} \frac{a_{1}+a_{2}+\cdots+a_{n}}{n}=0\) ，是否有 \(\lim _{n \rightarrow \infty} a_{n}=0\) ？为什么？
\end{itemize}

\subsection{题目5}
原题边界“（3）”：求 \(\lim _{n \rightarrow \infty} \sqrt[n]{\frac{1}{n!}}\) ．

\subsection{题目2}
二．（15 分）已知 \(f(x, y)\) 在 \((0,0)\) 的某邻域连续，且 \(F(t)=\iint_{x^{2}+y^{2} \leq t^{2}} f(x, y) \mathrm{d} x \mathrm{~d} y\) ，求 \(\lim _{t \rightarrow 0} \frac{F^{\prime}(t)}{e^{3 t}-1}\) ．

\subsection{题目3}
三。（15 分）解答如下问题：

\begin{itemize}
  \item[（1）] 证明：若 \(\lim _{n \rightarrow \infty}\left(a_{n+1}-a_{n}\right)=a\) ，则 \(\lim _{n \rightarrow \infty} \frac{a_{n}}{n}=a\) ．
\end{itemize}

\subsection{题目3}
原题边界“（2）”：已知 \(0<x_{1}<1, x_{n+1}=x_{n}\left(1-x_{n}\right)\) ，证明： \(\lim _{n \rightarrow \infty} n x_{n}=1\) ．

\subsection{题目4}
四．（15 分）设函数 \(f(x)\) 在 \(x=0\) 的某邻域存在连续的一阶导数，且 \(f^{\prime}(0)=0, f^{\prime \prime}(0)\) 存在，求极限

\[
\lim _{x \rightarrow 0} \frac{f(x)-f(\ln (1+x))}{x^{3}} .
\]

\subsection{题目5}
五．（15 分）求 \(\sum_{k=1}^{\infty}\left(\frac{(-1)^{k}}{2 k+1}+\frac{1}{k(k+1)}\right)\) ．

\subsection{题目6}
六．（15 分）求第二型曲面积分

\[
\begin{gathered}
\iint_{S}(y+z) \mathrm{d} y \mathrm{~d} z+\left(x^{2}+z^{2}\right) y \mathrm{~d} z \mathrm{~d} x+(x+y) \mathrm{d} x \mathrm{~d} y . \\
\left(\text { 另一个版本为 } \iint_{S} y z \mathrm{~d} y \mathrm{~d} z+\left(x^{2}+z^{2}\right) y \mathrm{~d} z \mathrm{~d} x+x y \mathrm{~d} x \mathrm{~d} y\right)
\end{gathered}
\]

其中 \(S: 4-y=x^{2}+z^{2}(y \geq 0)\) 的右侧。

\subsection{题目7}
七．（15 分）叙述并证明康托尔（Cantor）定理．

\subsection{题目8}
八．（15 分）设 \(f(x, y)=|x-y| \varphi(x, y)\) ，其中 \(\varphi(x, y)\) 在 \((0,0)\) 的某邻域内有定义，且在 \((0,0)\) 处连续．证明：\(\varphi(0,0)=0\) 的充分必要条件是 \(f(x, y)\) 在 \((0,0)\) 处可微．

\subsection{题目9}
九．（15 分）已知黎曼函数

\[
R(x)= \begin{cases}\frac{1}{q}, & x=\frac{p}{q}, \text { 其中 } p, q \text { 为正整数, 且 }(p, q)=1 ; \\ 0, & x \text { 为无理数 } ; \\ 1, & x=0,1 .\end{cases}
\]

讨论并说明 \(R(x)\) 在 \([0,1]\) 上的可积性与连续性．

附加题．（15 分）已知 \(g(0)=\frac{\sqrt{\pi}}{2}\) ，求 \(g(\alpha)=\int_{0}^{+\infty} e^{-x^{2}} \cos 2 \alpha x \mathrm{~d} x\) ．
\end{document}
